\subsection{Check 1: Viability}
Nearest hole to the pupil center: the hole at $(0,0)$ itself, $d=0$.
\begin{equation}
d - R_h = 0 - 1000 = -1000\mu\text{m} \;<\; R_p = 2000\mu\text{m} \quad\Rightarrow\quad \text{SAFE}
\end{equation}
This design has a hole sitting exactly on the optical axis --- viability is
not a close call here, unlike the two failure cases documented earlier in
this project where the nearest hole missed the pupil by $11.6\mu$m in
\emph{either} direction.

\subsection{Check 2: Ghost Margin}
Nearest hole pair (excluding the center hole, which has no close neighbor
in this direction): $(0,0)$ to $(2000,2000)$,
\begin{equation}
s = \sqrt{2000^2+2000^2} = 2828.4\mu\text{m}
\end{equation}
At $M=3$D, $P_0=60$D:
\begin{equation}
\Delta r = \frac{2828.4 \times 3}{60-3} = 148.9\mu\text{m}
\end{equation}
This must be compared against the resolvable-blur threshold for $M{=}3$D to
determine pass/fail --- the number itself is the part to hand-check; the
threshold comparison reuses whatever value the pipeline already reports for
that $M$.

\subsection{Check 3: Packing Density}
Convex hull area of the 9 points (shoelace formula on the outer square,
side $8000\mu$m, rotated $45^\circ$): $\text{Hull}_{\text{area}} =
3.200\times10^{7}\mu\text{m}^2$. Hull perimeter $\approx 22{,}627\mu$m.
\begin{equation}
\text{Area} = 3.200\times10^7 + 22{,}627\times1000 + \pi(1000)^2 = 5.777\times10^7\mu\text{m}^2
\end{equation}
\begin{equation}
\text{density} = \frac{9}{5.777\times10^7} = 1.558\times10^{-7}\;\mu\text{m}^{-2}
\end{equation}

\subsection{Cross-Check Against the Pipeline}
\begin{center}
\begin{tabular}{lrr}
\toprule
Quantity & Hand-computed & Pipeline reported \\
\midrule
Packing density ($\mu$m$^{-2}$) & $1.558\times10^{-7}$ & $1.5579\times10^{-7}$ \\
\bottomrule
\end{tabular}
\end{center}
Match to 4 significant figures. This is the standard the hand-check is held
to: agreement within rounding, not exact bit-for-bit equality.

\newpage
